% Unit 4 map -- one page.
\documentclass{apchem-worksheet}

\apcunit{4}
\apcunittitle{Chemical Reactions}
\apcdoctitle{Unit Map}
\apcsubtitle{CED Unit 4 \textbullet\ 7--9\% of the AP Exam \textbullet\ 7 blocks + exam}

\begin{document}
\apctitleblock

\begin{apcnote}[How this unit is graded --- and what makes it different]
Unit exam Mon Oct 26: 20 multiple choice (\textbf{no calculator}) + 1 long
free-response (10 points), 52 minutes.

If you worked Zumdahl Chapters 3 and 4, you already have most of the
\emph{mechanics} of this unit. What AP adds is the demand to
\textbf{represent} reactions --- particulate diagrams, net ionic equations,
and evidence-based claims about what actually happened. Expect to draw as
often as you calculate.
\end{apcnote}

\section*{\sffamily Block calendar}

\renewcommand{\arraystretch}{1.3}
\noindent\begin{tabular}{@{}p{2.1cm}p{5.0cm}p{1.5cm}p{3.9cm}p{1.9cm}@{}}
\toprule
\textbf{Date} & \textbf{Topics} & \textbf{CED} & \textbf{Read (PDF pp.)} & \textbf{Practice}\\
\midrule
Fri Oct 9 & Physical vs chemical change; introducing reactions & 4.1, 4.4 & \S1.10 \textbullet\ 53--57; \S3.8 \textbullet\ 132--134 & WS 1\\
Mon Oct 12 & Representations: symbolic and particulate & 4.3 & \S4.6--4.7 \textbullet\ 188--192 & WS 1\\
Wed Oct 14 & Net ionic equations & 4.2 & \S4.6 \textbullet\ 188--190 & WS 2\\
Fri Oct 16 & Types of reactions; precipitation & 4.7 & \S4.4--4.5 \textbullet\ 182--188 & WS 2\\
Mon Oct 19 & Acid--base reactions & 4.8 & \S4.8 \textbullet\ 192--199 & WS 2\\
Wed Oct 21 & Oxidation--reduction reactions & 4.9 & \S4.9--4.10 \textbullet\ 199--210 & WS 3\\
Fri Oct 23 & Stoichiometry and titration & 4.5, 4.6 & \S3.9--3.11 \textbullet\ 134--153; \S4.11 \textbullet\ 210--212 & WS 4 (FRQ)\\
\textbf{Mon Oct 26} & \textbf{UNIT 4 EXAM} & all & review sheet & ---\\
\bottomrule
\end{tabular}

{\sffamily\footnotesize Dates assume continuous instructional weeks; shift
for any co-op breaks.}

\vspace{0.5em}
\section*{\sffamily By the end of this unit you can\ldots}

\begin{itemize}[leftmargin=*,itemsep=0.12em]
  \item Distinguish physical from chemical change using particle-level
        evidence, not just observations. \ced{4.4}
  \item Translate among symbolic, particulate, and macroscopic
        representations of the same reaction. \ced{4.3}
  \item Write balanced formula, complete ionic, and net ionic equations, and
        justify which species are split. \ced{4.2}
  \item Classify a reaction as precipitation, acid--base, or redox, and
        predict its products. \ced{4.7}
  \item Write net ionic equations for strong and weak acid neutralizations.
        \ced{4.8}
  \item Assign oxidation states, identify what is oxidized and reduced, and
        name both agents. \ced{4.9}
  \item Carry out solution stoichiometry, including limiting reactant.
        \ced{4.5}
  \item Analyze a titration: equivalence point, curve interpretation, and
        concentration determination. \ced{4.6}
\end{itemize}

\begin{apctrap}[Where students lose points in this unit]
Splitting weak acids or water in a net ionic equation; drawing particulate
diagrams that violate conservation of atoms; calling a colour change
``proof'' of chemical reaction without ruling out dilution or indicator
effects; naming the oxidizing agent when they mean the reducing agent; and
using a 1:1 titration ratio for a diprotic acid.
\end{apctrap}

\end{document}
